\chapter*{引言}
\addcontentsline{toc}{chapter}{引言}
\VerdierPhysicalPageAnchor{0018}
\VerdierSemanticAnchor{intro}

\VerdierSemanticAnchor{intro-opening}
我们在本文中提出一种超同调的形式体系。

\VerdierSemanticAnchor{intro-1}
\noindent\textbf{1.}\quad
设 $\mathcal A$ 与 $\mathcal B$ 是两个阿贝尔范畴，
$F:\mathcal A\to\mathcal B$ 是加性函子。当 $\mathcal A$ 有足够多的
内射对象时，$F$ 的右导出函子的构造如今已是经典的 [1]：
$\mathcal A$ 的每个对象 $X$ 都有一个内射分解
$X\xrightarrow{\varepsilon}I^{\bullet}(X)$，且它在“同伦意义下”唯一。
因此，复形 $F(I^{\bullet}(X))$ 的上同调对象由 $X$ 在典范同构意义下
确定，并且函子性地依赖于 $X$。由此得到的函子就是 $F$ 的右导出函子。

\VerdierSemanticAnchor{intro-1-homotopy-quasi}
这样，我们把一个在“同伦意义下”确定的 $\mathcal B$ 中对象复形函子性地
对应给 $\mathcal A$ 的每个对象 $X$，然后取其上同调。然而，同调代数中的
通常构造往往并不给出在同伦意义下确定的复形，而只给出在
\emph{拟同构意义下}确定的复形（如果一个复形态射在上同调对象上诱导同构，
就称它为拟同构）。下面用两个例子说明这一点。

\VerdierSemanticAnchor{intro-1-f-acyclic}
首先，为了研究 $F$ 的导出函子，应当考虑比内射分解更一般的分解
$X\xrightarrow{\varepsilon'}\mathcal R(X)$：即由对函子 $F$ 无环的对象
所组成的分解（在层论中，是由松弛层或软层组成的分解 [23]；在模论中，
是由平坦模组成的分解，等等）。这样的分解在同伦意义下不再必然唯一。
但是，给定一个 $F$-无环分解
$X\xrightarrow{\varepsilon'}\mathcal R(X)$ 和一个内射分解
$X\xrightarrow{\varepsilon}I^{\bullet}(X)$，存在唯一的（在同伦意义下）
分解态射，即一个复形态射
$\rho:\mathcal R(X)\to I^{\bullet}(X)$，使下页中的图交换：
